\section{Existing project bounds by scope}
\label{sec:ledger}

\begin{center}
\small
\begin{tabularx}{\textwidth}{@{}p{0.25\textwidth}p{0.25\textwidth}X@{}}
\toprule
Result & Actual scope & What it does not prove \\
\midrule
APPR star lower bound
& Every legal active ordering of the stated APPR push rule
& A lower bound for arbitrary PageRank algorithms or SDD response \\
Coordinate RPPR star lower bound
& Thresholded coordinate ISTA and its stated coarse-to-fine wrapper
& A lower bound for accelerated, block, or elimination methods \\
CF-Push/SOR spider bound
& The fixed-relaxation FIFO tail after the stated coarse gate
& A lower bound for variable polynomial relaxation or response deflation \\
AESP--LocGD star bound
& The literal AESP outer loop with batched LocGD inner solves
& A lower bound for arbitrary AESP inner maps \\
ASPR path bound
& Literal repeated principal APGD calls and fresh discovery scans
& A lower bound for no-restart continuation or implicit response \\
Persistent-support path bundle
& The note's repeated one-hop resident-volume oracle
& A lower bound once directed elimination messages are allowed \\
Path materialization barrier
& Explicit active-vector output after every nested expansion
& A barrier to an implicit append-only factor and final reverse solve \\
\bottomrule
\end{tabularx}
\end{center}

This ledger explains why the same graph can be hard for one local solver and
easy for another.  The hard resource may be repeated mass processing,
resident-volume rescanning, loss of conjugacy, reflection, factor
materialization, or boundary reporting.  These mechanisms require different
model clauses.
